\documentclass[
  a4paper, 11pt, fleqn, twoside,
  abstract=true
]{scrartcl}
\usepackage[ngerman]{babel}
\usepackage[T1]{fontenc}

\newif\ifscreen
\screentrue

\usepackage{xcolor}
\usepackage{amsmath}
\usepackage{amsthm}
% \usepackage{mathtools}

\usepackage{libertine}
\usepackage[libertine]{newtxmath}
\usepackage[scaled=0.78]{DejaVuSansMono}
% \usepackage{bm}

% \usepackage{tikz}
% \usepackage{tikz-cd}
\usepackage{booktabs}
\usepackage{listings}
\usepackage{microtype}
% \usepackage{lipsum}

% \usepackage{enumitem}
% \setlist{nosep}

\ifscreen
  \definecolor{clink}{RGB}{0, 60, 140}
  \definecolor{ccite}{RGB}{120, 0, 160}
  \definecolor{curl}{RGB}{180, 0, 180}
\else
  \definecolor{clink}{RGB}{0, 0, 0}
  \definecolor{ccite}{RGB}{0, 0, 0}
  \definecolor{curl}{RGB}{0, 0, 0}
\fi

\usepackage{hyperref}
\hypersetup{
  colorlinks=true,
  linkcolor=clink,
  citecolor=ccite,
  urlcolor=curl,
  pdftitle={Ententheorie},
  pdfauthor={Donald Duck}
}

\usepackage[a4paper,
  left=30mm, right=30mm, top=16mm, bottom=30mm,
  includehead, headheight=5mm
]{geometry}

\usepackage[headsepline=0.8pt]{scrlayer-scrpage}
\clearpairofpagestyles
\automark[subsection]{section}
\ohead{\pagemark}
\ihead{\headmark}
\addtokomafont{pagehead}{\upshape}

% \linespread{1.04}
\setkomafont{disposition}{\normalfont\bfseries}

\makeatletter
\renewenvironment{proof}[1][\proofname]{\par
  \normalfont\topsep2\p@\@plus2\p@\relax
  \trivlist
  \item[\hskip\labelsep\bfseries #1\@addpunct{.}]\ignorespaces
}{\endtrivlist\@endpefalse}
\makeatother

\theoremstyle{plain}
\newtheorem{theorem}{Satz}[section]
% \newtheorem{lemma}[theorem]{Lemma}
% \newtheorem{corollary}[theorem]{Korollar}
% \newtheorem{proposition}[theorem]{Proposition}

\theoremstyle{definition}
\newtheorem{definition}[theorem]{Definition}
% \newtheorem{example}[theorem]{Beispiel}

\newcommand{\strong}[1]{\textbf{#1}}
\newcommand{\keywords}[1]{%
  \vspace{0.6em}\noindent\strong{Schlagwörter:} #1}

\lstset{
  basicstyle=\ttfamily,
  numbers=left,
  xleftmargin=\mathindent,
  % showstringspaces=false,
  upquote=true
}

\newcommand{\N}{\mathbb N}
\newcommand{\Z}{\mathbb Z}
\newcommand{\Q}{\mathbb Q}
\newcommand{\R}{\mathbb R}
\newcommand{\C}{\mathbb C}
% \newcommand{\ui}{\mathrm i}
% \newcommand{\ee}{\mathrm e}

% \newcommand{\cond}{\rightarrow}
% \newcommand{\bicond}{\leftrightarrow}
% \newcommand{\lnec}{\Box}
% \newcommand{\lpos}{\Diamond}

% \newcommand{\compc}{\mathsf c}
% \newcommand{\setsystem}[1]{\mathcal{#1}}
% \newcommand{\powerset}{\setsystem P}
% \newcommand{\category}[1]{\mathcal{#1}}
% \DeclareMathOperator{\id}{id}

\title{Ententheorie}
\subtitle{Eine konservative Vorlage für\\ mathematische Artikel}
\author{Donald Duck\\ \large\texttt{duck@example.com}}
\date{1. Jan. 2000}

\begin{document}

\maketitle

\begin{abstract}\noindent
Es wird der Kontraktionssatz zur Fixpunktiteration gezeigt und damit
anschließend als praktisches Beispiel ein Berechnungsverfahren
diskutiert.

\keywords{Fixpunktiteration, Fixpunktsatz von Banach,
Kontraktion, Grenzwert, Dehnungsbeschränktheit (Lipschitz-Stetigkeit),
Cauchyfolge, geometrische Summenformel.}
\end{abstract}

\tableofcontents

\section{Theoretische Überlegungen}

\subsection{Fixpunktiteration}

Zu einer in Fixpunktform $x=f(x)$ gebrachten Gleichung lassen sich
Lösungen unter Umständen mit der Iteration $x_{n+1} := f(x_n)$, und
dies dann recht unempfindlich vom Startwert $x_0$, beliebig gut annähern.
Im Folgenden wird untersucht, unter welchen Umständen sich die Iteration
gutmütig verhalten muss. Wir folgen den Ausführungen in \cite{Heuser1}.

\begin{theorem}
Zu einer monoton steigenden stetigen Funktion $f\colon [a,b]\to [a,b]$
konvergiert die Folge $x_{n+1} := f(x_n)$ zu jedem $x_0\in [a,b]$ monoton
gegen einen Fixpunkt von $f$.
\end{theorem}

\begin{proof}
Induktiv klärt sich, dass $x_n\in [a,b]$ und $x_{n+1}\ge x_n$ bei
$x_1\ge x_0$ oder $x_{n+1}\le x_n$ bei $x_1\le x_0$ für jedes
$n\in\N$ gilt. Da die Folge $(x_n)$ also in jedem Fall monoton und durch
$a,b$ beschränkt ist, muss sie einen Grenzwert $c\in [a,b]$ besitzen.
Und weil $f$ stetig ist, gilt zu diesem die Termumformung
\[f(c) = f(\lim_{n\to\infty} x_n) = \lim_{n\to\infty} f(x_n) =
\lim_{n\to\infty} x_{n+1} = \lim_{n\to\infty} x_n = c,\]
kurz $f(c) = c$, womit $c$ ein Fixpunkt von $f$ sein muss.\,\qedsymbol
\end{proof}

\subsection{Kontraktionssatz}

Der folgende Spezialfall des banachschen Fixpunktsatzes zeigt Umstände
auf, unter denen der Fixpunkt eindeutig bestimmt ist.

\begin{definition}[Dehnungsbeschränktheit]
Eine Funktion $f\colon X\to\R$ mit $X\subseteq\R$ heißt
\emph{dehnungsbeschränkt} oder \emph{Lipschitz"=stetig}, wenn es eine
Konstante $q\ge 0$ gibt, so dass $|f(x)-f(y)| \le q|x-y|$ für alle
$x,y\in X$ gilt.
\end{definition}

\noindent
Grafisch gesehen wird ein Doppelkegel mit festem Anstieg $q$ an einen
beliebigen Punkt auf dem Funktionsgraphen angeheftet, so dass die
Kegelspitzen auf dem Graph liegen. Der Funktionsgraph darf den Kegel
berühren, aber nicht in ihn eintauchen.

\begin{theorem}[Kontraktionssatz]
Eine Funktion $f\colon [a,b]\to [a,b]$, die mit Dehnungsschranke $q<1$
kontrahiert, besitzt genau einen Fixpunkt $c$. Dieser ist Grenzwert
der Folge $x_{n+1} := f(x_n)$, wobei der Startwert $x_0\in [a,b]$ frei
gewählt werden darf. Hierbei gilt die Fehlerabschätzung
\[|c-x_n| \le \frac{q^n}{1-q}|x_1-x_0|.\]
\end{theorem}

\begin{proof}
Laut der Kontraktionseigenschaft gilt
\[|x_{n+1}-x_n| = |f(x_n)-f(x_{n-1})| \le q|x_n-x_{n-1}|\]
für $n\ge 1$. Induktiv erhält man somit $|x_{n+1}-x_n|\le q^n|x_1-x_0|$.
Via Teleskopieren, Dreiecksungleichung und geometrischer Summenformel
findet sich zu jedem $k$ demnach
\begin{align*}
|x_{n+k}-x_n| &= \bigg|\sum_{i=1}^k (x_{n+i}-x_{n+i-1})\bigg|
\le \sum_{i=1}^k |x_{n+i}-x_{n+i-1}|\\
&\le \sum_{i=1}^k q^{n+i-1}|x_1-x_0|
= q^n\frac{1-q^k}{1-q}|x_1-x_0|\le \frac{q^n}{1-q}|x_1-x_0|.
\end{align*}
Insofern $q<1$ ist, muss $(x_n)$ somit eine Cauchyfolge mit Grenzwert
$c\in [a,b]$ sein. Weil $f$ als dehnungsbeschränkte Funktion erst
recht stetig ist, hat man außerdem wieder $f(c)=c$, womit $c$ ein
Fixpunkt von $f$ ist. Ist $d$ mit $f(d)=d$ nun ein beliebiger weiterer
Fixpunkt, dann gilt aufgrund der Kontraktionseigenschaft die Ungleichung
\[|c-d| = |f(c)-f(d)|\le q|c-d|,\]
was aber nur für $|c-d|=0$, also $c=d$ möglich ist. Demnach ist der
Fixpunkt eindeutig. Die Fehlerabschätzung ergibt sich schließlich bei
$k\to\infty$.\,\qedsymbol
\end{proof}

\section{Eine Beispieldiskussion}

\subsection{Aufstellung der Iteration}

Zur praktischen Anwendung des Kontraktionssatzes wollen wir ein
Verfahren zur Berechnung der Quadratwurzel aufstellen, wenngleich
es mit dem Heron"= bzw. Newton"=Verfahren wesentlich
effizientere Ansätze geben mag.

Zur Berechnung der Quadratwurzel $\sqrt{a}$ bietet sich
die Iteration $x_{n+1} := f(x_n)$ mit
\[X :=[\sqrt{a},a+1],\quad f\colon X\to X, \quad
f(x) := x + \frac{a-x^2}{2(a+1)}\]
an. Monotones Steigen liegt vor, falls $f'(x)\ge 0$ gilt. Mit
$f'(x) = 1 - \frac{x}{a+1}$ führt die Forderung auf die Ungleichung
$\frac{x}{a+1}\le 1$, die wegen $a>0$ äquivalent zu $x\le a+1$ umgeformt
werden darf. Der Graph der Funktion $f'$ ist genau die von $(0,1)$
nach $(a+1,0)$ verlaufende Gerade. Als kleinstmögliche Dehnungsschranke
bekommt man daher $q=f'(\sqrt{a})$.\footnote{Streng genommen wäre hier
explizit der Mittelwertsatz der Differentialrechnung anzuwenden, um den
Bezug zwischen maximalem Anstieg und Dehnungsbeschränktheit herzustellen:
Laut diesem gibt es zu beliebigen $x,y\in X$ ein $u\in X$, so dass
$|f(x)-f(y)| = |f'(u)||x-y|$. Weiter gilt immer $|f'(u)|\le q$, insofern
$q$ der maximale Anstieg ist. Demnach gilt $|f(x)-f(y)| \le q|x-y|$.}

Man überzeugt sich unschwer von $\sqrt{a}<a+1$, womit das Intervall $X$
wohldefiniert ist. Es verbleibt noch zu prüfen, ob $f$ mit $f(X)\subseteq X$
wohldefiniert ist. Da $f$ monoton steigt, genügt es,
$[f(\sqrt{a}),f(a+1)]\subseteq X$ zu prüfen, also $\sqrt{a}\le f(\sqrt{a})$
und $f(a+1)\le a+1$, was unschwer zu bewältigen ist.

\subsection{Berechnung der Iteration}

Wir führen zu $a := 2$ eine Beispielrechnung durch. Man hat
$q\approx 0.528595$ als Dehnungsschranke, was bedeutet, dass sich der
abgeschätzte Fehler in jedem Schritt ungefähr halbiert. Die Berechnung
wird gemäß Listing \ref{lst:Berechnung} ausgeführt und ist in
Tabelle \ref{tab:Berechnung} aufgeführt.

\begin{table}
\begin{center}
\caption{Berechnung von $\sqrt{2}$}\label{tab:Berechnung}
\begin{tabular}{@{}cc@{}}
\begin{tabular}{rccc}
\toprule
$n$ & $x_n$ & $|c-x_n|$ & $\frac{q^n}{1-q}|x_1-x_0|$\\
\midrule
$ 1$ & $1.833333$ & $4.19\times 10^{-1}$ & $1.31\times 10^{-0}$\\
$ 2$ & $1.606481$ & $1.92\times 10^{-1}$ & $6.92\times 10^{-1}$\\
$ 3$ & $1.509684$ & $9.55\times 10^{-2}$ & $3.66\times 10^{-1}$\\
$ 4$ & $1.463160$ & $4.89\times 10^{-2}$ & $1.93\times 10^{-1}$\\
\midrule
$ 5$ & $1.439687$ & $2.55\times 10^{-2}$ & $1.02\times 10^{-1}$\\
$ 6$ & $1.427571$ & $1.34\times 10^{-2}$ & $5.40\times 10^{-2}$\\
$ 7$ & $1.421244$ & $7.03\times 10^{-3}$ & $2.85\times 10^{-2}$\\
$ 8$ & $1.417922$ & $3.71\times 10^{-3}$ & $1.51\times 10^{-2}$\\
\midrule
$ 9$ & $1.416171$ & $1.96\times 10^{-3}$ & $7.97\times 10^{-3}$\\
$10$ & $1.415248$ & $1.03\times 10^{-3}$ & $4.21\times 10^{-3}$\\
$11$ & $1.414760$ & $5.47\times 10^{-4}$ & $2.23\times 10^{-3}$\\
$12$ & $1.414502$ & $2.89\times 10^{-4}$ & $1.18\times 10^{-3}$\\
\bottomrule
\end{tabular}
&
\begin{tabular}{rccc}
\toprule
$n$ & $x_n$ & $|c-x_n|$ & $\frac{q^n}{1-q}|x_1-x_0|$\\
\midrule
$13$ & $1.414366$ & $1.53\times 10^{-4}$ & $6.23\times 10^{-4}$\\
$14$ & $1.414294$ & $8.07\times 10^{-5}$ & $3.29\times 10^{-4}$\\
$15$ & $1.414256$ & $4.27\times 10^{-5}$ & $1.74\times 10^{-4}$\\
$16$ & $1.414236$ & $2.25\times 10^{-5}$ & $9.19\times 10^{-5}$\\
\midrule
$17$ & $1.414225$ & $1.19\times 10^{-5}$ & $4.86\times 10^{-5}$\\
$18$ & $1.414220$ & $6.30\times 10^{-6}$ & $2.57\times 10^{-5}$\\
$19$ & $1.414217$ & $3.33\times 10^{-6}$ & $1.36\times 10^{-5}$\\
$20$ & $1.414215$ & $1.76\times 10^{-6}$ & $7.18\times 10^{-6}$\\
\midrule
$21$ & $1.414214$ & $9.30\times 10^{-7}$ & $3.79\times 10^{-6}$\\
$22$ & $1.414214$ & $4.92\times 10^{-7}$ & $2.01\times 10^{-6}$\\
$23$ & $1.414214$ & $2.60\times 10^{-7}$ & $1.06\times 10^{-6}$\\
$24$ & $1.414214$ & $1.37\times 10^{-7}$ & $5.60\times 10^{-7}$\\
\bottomrule
\end{tabular}
\end{tabular}
\end{center}
\end{table}

\begin{figure}
\begin{lstlisting}[language=Python,
caption={Berechnungsverfahren}, label={lst:Berechnung}]
from math import sqrt

a = 2
q = 1 - sqrt(a)/(a + 1)
x0 = a + 1

def f(x):
    return x + (a - x*x)/(2*a + 2)

def err_esti(n):
    return q**n/(1 - q)*abs(f(x0) - x0)

x = x0; c = sqrt(a)
for n in range(1, 25):
    x = f(x)
    e1 = abs(c - x); e2 = err_esti(n)
    print(f"{n:2}  {x:.6f}  {e1:.2E}  {e2:.2E}")
    if n%4 == 0: print()
\end{lstlisting}
\end{figure}

% \lipsum[1-40]

\begin{thebibliography}{00}
\bibitem{Heuser1} Harro Heuser: \emph{Lehrbuch der Analysis. Teil 1}.
  Vieweg+Teubner, 17. Auflage, Wiesbaden 2009.\\
  \href{https://doi.org/10.1007/978-3-322-94702-4}{doi:10.1007/978-3-322-94702-4}.
  ISBN 978-3-8348-0777-9.
\end{thebibliography}

\end{document}
